\documentclass[level=2]{kidscode}
\begin{document}
\kctitle{Notes for Families}{What your child will learn in Level 2, how we keep them safe, and how you can help}

\section{Dear families}
Welcome to Level 2 of our home coding club! This letter is for every family whose child takes part,
including friends who join us from another home. You do not need to know anything about
computers to help: the most useful thing you can do is be curious about what your child made.

In Level 1 the children moved from Scratch blocks to typing real Python code. In Level 2 they
become real programmers. Over 12 weekly sessions of one hour they will learn to organise programs
into functions, keep collections of things in lists and dictionaries, save high scores in files,
write automatic tests that check their code, and build arcade games that react to the keyboard.
It ends with a final project, an arcade game or an art studio of their own design, and a
showcase where you are warmly invited to play.

\begin{tabularx}{\linewidth}{@{}>{\bfseries}p{4.2cm}X@{}}
\toprule
When & Once a week, 60 minutes, for about 14--16 weeks (12 sessions plus 2 spare sessions for holidays). \\
Where & At our home, on our Macs, using the free program Thonny. \\
Group & Up to 4 children, working in pairs. \\
Homework & 15--20 minutes a week. \enquote{Core} is for everyone; \enquote{Challenge} and \enquote{Creative} are optional. \\
Cost & Nothing to buy. Everything is printed or on our computers. \\
Showcase & Week 12. Families are invited to play the games and see the art. Date: \rule{4cm}{0.4pt} \\
\bottomrule
\end{tabularx}

\section{Our safety and privacy promises}
\begin{safety}[What we promise]
\begin{itemize}[itemsep=2pt]
  \item \textbf{No accounts, no logins, no uploading.} Everything stays on the computers in our home.
        No child needs an e-mail address or any online account.
  \item \textbf{No personal information in programs.} Children use a nickname or \enquote{hero name}.
        Programs never ask for real names, ages, addresses, schools, photos or passwords. We teach
        why in week 1, and come back to it all term.
  \item \textbf{No photos of children} are taken or shared without your permission. If you allow
        photos at the showcase, we photograph screens and projects, not faces or name labels.
  \item \textbf{Credit and kindness.} Children learn to credit other people's art and code, give
        kind feedback (\enquote{two stars and a wish}) and design games that are fair to players.
  \item \textbf{Screen balance.} Sessions include unplugged (paper and movement) activities and a
        stretch break. Homework is short.
  \item \textbf{Accessible for everyone.} Colours are chosen to be safe for colour-blind players,
        and a dyslexia-friendly edition of every printed page is available: just ask.
\end{itemize}
\end{safety}

\section{Homework at home}
Most homework can be done on paper. Some tasks ask the child to change a program. If your child
does not have a computer with Thonny at home, that is completely fine: they can do the paper parts
and finish the rest at the start of the next session. If you would like to install it, Thonny is
free for Mac, Windows and Linux (\texttt{thonny.org}); we can copy the course folder onto a USB
stick for you. No account is needed.

\section{Ask your child about\ldots}
A good question at dinner does more than any worksheet. Here is one for each week.

{\small
\renewcommand{\arraystretch}{1.3}
\begin{longtable}{@{}>{\bfseries}p{0.8cm}>{\raggedright\arraybackslash}p{4.0cm}>{\raggedright\arraybackslash}p{9.8cm}@{}}
\toprule
Wk & Topic & Ask your child\ldots \\ \midrule\endhead
1 & Types and messages & \enquote{Why does Python get confused by \texttt{"5" + 5}?} and \enquote{What is your hero name, and why do we not use real names?} \\
2 & Functions that give back & \enquote{What is the difference between a function that \emph{prints} and one that \emph{returns}?} \\
3 & Lists & \enquote{Why does a list start counting at 0?} Ask them to show you a colour palette made from a list. \\
4 & While loops & \enquote{What is an infinite loop, and how do you stop one?} Play their guess-my-number game. \\
5 & Dictionaries and secret codes & Ask them to send you a message in Caesar cipher, and decode one you write. \enquote{Why should passwords be secret?} \\
6 & Review and checkpoint & \enquote{What do you do when you get an error message?} (Hint: read the last line, print things, or explain it to a rubber duck.) \\
7 & Files & \enquote{How does a game remember your high score after you close it?} \enquote{What should a program \emph{not} save?} \\
8 & Testing & \enquote{What is an edge case?} Ask them to think of tests for the toaster or the microwave. \\
9 & Games with keys and timers & \enquote{How does the game know when two things touch?} Ask about their final project idea. \\
10 & Project studio 1 & \enquote{What is the first thing your project has to do? How did you split the work with your partner?} \\
11 & Project studio 2 & \enquote{What did a playtester tell you, and what did you change?} \\
12 & Showcase & Come and play! Ask: \enquote{Which part was hardest, and how did you solve it?} \\
\bottomrule
\end{longtable}}

\section{How you can help}
\begin{itemize}[itemsep=2pt]
  \item \textbf{Be the audience.} Ask to see what they made this week. Play their games, even
        the very hard or very silly ones.
  \item \textbf{Be the rubber duck.} When something does not work, ask them to explain the code to
        you line by line. Often they find the bug while explaining. You do not need to understand
        the code to help this way.
  \item \textbf{Celebrate bugs.} Mistakes are a normal part of programming. Praise the effort and
        the strategy (\enquote{you read the error message carefully!}) more than the result.
  \item \textbf{Please do not fix the code for them.} Questions help more than answers:
        \enquote{What did you expect to happen? What happened instead?}
  \item \textbf{Keep it light.} Homework should take 15--20 minutes. If it takes much longer,
        let them stop and let us know.
  \item \textbf{Keep their files in one folder,} with clear names, so nothing gets lost.
\end{itemize}

\begin{tip}[Questions or concerns?]
Please talk to us at any time, for example about learning needs, the dyslexia-friendly edition,
photos, or days your child cannot come.\\[6pt]
Contact: \rule{7cm}{0.4pt}
\end{tip}

\end{document}
